\documentclass[11pt]{article}
\usepackage[a4paper,margin=24mm]{geometry}
\usepackage[T1]{fontenc}
\usepackage{lmodern}
\usepackage{microtype}
\usepackage{graphicx}
\usepackage{xcolor}
\usepackage{caption}
\usepackage{enumitem}
\usepackage{float}
\usepackage{fancyhdr}
\usepackage[hidelinks]{hyperref}

\definecolor{Accent}{HTML}{315B7D}
\definecolor{Muted}{HTML}{5B6570}

\captionsetup{font=small,labelfont={bf,color=Accent}}
\setlist[itemize]{leftmargin=1.3em,itemsep=2pt,topsep=3pt}
\newcommand{\Faq}{\textit{F.~aquilonia}}
\newcommand{\Fpo}{\textit{F.~polyctena}}
\newcommand{\pt}[1]{\medskip\noindent\textbf{\color{Accent}#1}\par\nobreak\smallskip}

\pagestyle{fancy}\fancyhf{}
\lhead{\small\color{Muted}Ancestry sorting at diagnostic loci}
\rhead{\small\color{Muted}talk notes}
\cfoot{\small\thepage}
\renewcommand{\headrulewidth}{0.4pt}
\setlength{\parskip}{0.4em}\setlength{\parindent}{0pt}

\title{\vspace{-3.2em}\textbf{\color{Accent}Ancestry sorting at the diagnostic loci}\\[0.2em]
\large What the hybrids keep, and where --- a walk through the figures}
\author{\normalsize Poikela et al. \; -- \; talk-preparation summary}
\date{\normalsize\today}

\begin{document}
\maketitle
\vspace{-1.8em}

\noindent\textit{These are notes to talk from, not a methods section. The thread:
we take the ancestry-diagnostic markers of the \Faq{}\,$\times$\,\Fpo{} hybrids,
reduce them to (nearly) independent units, and ask how much ancestry has sorted,
which way, and where in the genome.}

%% ------------------------------------------------------------------
\pt{The setup in one breath}
\begin{itemize}
  \item Twenty wild hybrid \textit{Formica} populations between \Faq{} and \Fpo{},
        spanning $\sim$2{,}000\,km.
  \item We focus on the \textbf{diagnostic markers} --- the $\sim$51{,}600 sites that
        tell the two parental species apart (diagnostic index $>-25$). These are the
        loci where ``which ancestry won'' is even a meaningful question.
  \item Colour code throughout: \textbf{purple = \Faq{}}, \textbf{yellow = \Fpo{}},
        teal = heterozygous / mixed.
\end{itemize}

%% ------------------------------------------------------------------
\pt{1.\; From 51{,}600 SNPs to 11{,}000 units --- why we don't count SNPs}
Linked markers tell the same story many times over. In the low-recombination
blocks a single ancestry pattern can be echoed by thousands of SNPs, so counting
SNPs counts the same thing repeatedly. We therefore rebuilt the linkage-disequilibrium
clustering \emph{from scratch, on the diagnostic markers only}, collapsing the
51{,}612 SNPs into \textbf{11{,}052 LD-reduced units} (a 5\,cM merge; $\sim$78\% reduction).
A caveat we keep coming back to: these units are \emph{much} more independent than
SNPs, but not perfectly independent.

\begin{figure}[H]\centering
\includegraphics[width=\textwidth]{../Figures/diem_circos_compare_snp_vs_ldreduced.png}
\caption{The same diagnostic data, same 195 individuals, same ordering (inner rings
= most \Faq{}). \textbf{Left:} every SNP --- speckly, with fat vertical blocks where
linked markers repeat. \textbf{Right:} the LD-reduced units --- the blocks collapse,
the picture thins by $\sim$5$\times$, the ancestry structure survives. This is the
whole reason we work with units, not SNPs.}
\end{figure}

%% ------------------------------------------------------------------
\pt{2.\; How much has sorted, and which way?}
A locus is ``sorted'' when the same ancestry has repeatedly become near-fixed across
independent hybrid populations. Two headlines, and they are exactly the kind of thing
that only shows up once you stop counting SNPs:

\begin{itemize}
  \item \textbf{LD reduction deflates the amount of sorting.} At our primary threshold,
        $\sim$13\% of diagnostic SNPs look ``sorted'' but only $\sim$7\% of independent
        units are --- the SNP number was inflated by linked repeats.
  \item \textbf{LD reduction can even flip the apparent \emph{direction}.} Among the most
        strongly sorted loci, the raw SNPs lean toward \Fpo{}; the independent units lean
        strongly toward \Faq{} ($\sim$80\%). The \Fpo{} ``signal'' was a handful of big
        low-recombination blocks, counted thousands of times. Collapse them and \Faq{}
        prevails.
\end{itemize}

\begin{figure}[H]\centering
\includegraphics[width=0.82\textwidth]{../Figures/diem_circos_sorting_sweep_snp.png}
\caption{Where the sorting is, and how it thins as we demand a stricter, more repeatable
signal (rings inner\,$\rightarrow$\,outer = increasingly strict). Purple = sorted toward
\Faq{}, yellow = toward \Fpo{}, grey = not sorted. The persistent \textbf{yellow blocks}
(e.g.\ chromosomes 5, 25, 26) are the large linkage blocks that make the per-SNP counts
lean \Fpo{} --- they are single units, not many.}
\end{figure}

%% ------------------------------------------------------------------
\pt{3.\; The same story, population by population}
Collapsing the 195 individuals into the 20 populations shows the diagnostic landscape
as a per-population map. The strongly sorted regions read as broad bands that are the
\emph{same colour across all populations} --- that is what ``repeatable sorting'' looks
like. Most of the genome is teal: these are balanced-ancestry hybrids, with sorting
concentrated at particular loci.

\begin{figure}[H]\centering
\includegraphics[width=0.72\textwidth]{../Figures/diem_circos_population_snp.png}
\caption{Twenty population rings, ordered by overall ancestry (\r{A}land most \Faq{},
Hein\"am\"aki most \Fpo{}). Each cell is the population's mean \Faq{}-allele frequency.
Broad purple/yellow bands = regions where the same ancestry has near-fixed across many
populations.}
\end{figure}

%% ------------------------------------------------------------------
\pt{4.\; Is sorting stronger where recombination is low?}
This is where we have to be careful --- and honest --- about effect sizes.

\begin{itemize}
  \item \textbf{Direction of the effect:} yes, among diagnostic loci sorting is weakly
        \emph{stronger} in low-recombination regions. Interestingly this is the
        \emph{opposite} of the genome-wide picture, because genome-wide the low-recombination
        regions are dominated by non-diagnostic markers that never sort. Condition on the
        diagnostic loci and the sign flips.
  \item \textbf{But the effect is small.} Recombination explains under $\sim$1.5\% of whether
        a unit sorts. The probability of sorting drifts from roughly 15\% down to 5\% across
        the well-sampled recombination range --- real and repeatable, but a gentle slope on
        a low baseline. The tiny $p$-values come from having thousands of units, not from a
        big effect.
  \item \textbf{Robust to how we test it.} The slope is negative at every sorting threshold,
        and it survives a genomic block bootstrap (resampling whole chromosomes) that respects
        the fact that units aren't perfectly independent.
  \item \textbf{The honest asterisk:} the very lowest-recombination regions can't really be
        judged at all --- their $\sim$10{,}000 diagnostic SNPs collapse to only $\sim$57
        independent units, so those points carry huge uncertainty (big error bars, small bars).
\end{itemize}

\begin{figure}[H]\centering
\includegraphics[width=\textwidth]{../Figures/di25_recomb_tau_sweep.png}
\caption{\textbf{Left:} fraction of units sorted vs recombination, at each sorting threshold
$\tau$; grey background bars = number of \emph{independent} units per decile (near zero on
the left --- that is the honest caveat). \textbf{Right:} the recombination effect at each
$\tau$ with block-bootstrap 95\% intervals --- all negative, all excluding zero, all small.}
\end{figure}

%% ------------------------------------------------------------------
\pt{Take-homes for the slide}
\begin{itemize}
  \item Ancestry sorting at diagnostic loci is \textbf{real and repeatable}, but you only
        measure it correctly once you \textbf{collapse linked markers into units} --- SNP
        counts both inflate the amount and can misstate the direction.
  \item Once cleaned up, the diagnostic loci sort \textbf{predominantly toward \Faq{}}.
  \item Sorting is \textbf{modestly concentrated in low-recombination regions} among diagnostic
        loci --- statistically solid, but a \textbf{small} effect, and the most extreme
        low-recombination regions are essentially single blocks we can't resolve.
  \item Throughout, the discipline is the same: \textbf{count independent units, report effect
        sizes, and let the significance follow} --- not the other way round.
\end{itemize}

\end{document}
